\documentclass[twoside]{article}
\usepackage[table]{xcolor}
\usepackage[paper=letterpaper,left=3cm,right=3cm,top=3cm,bottom=2cm,includefoot]{geometry}
\usepackage{fancyhdr}
\usepackage{tabularx}
\usepackage{multirow}
\usepackage[colorlinks]{hyperref}
\usepackage{hhline}
\usepackage{amsmath}
\usepackage{enumitem}
\setenumerate{itemsep=-3pt,topsep=3pt}
\setdescription{itemsep=-3pt,topsep=3pt,leftmargin=!,labelwidth=4.5cm}
\usepackage{pgfgantt}
\usepackage{graphicx}   \graphicspath{{./img/} {./tex/img/}}
\usepackage{amsfonts}
\usepackage[spanish]{babel}
\usepackage[utf8]{inputenc}
\usepackage[T1]{fontenc}
\usepackage{listings}
\usepackage{xcolor}
\usepackage[most]{tcolorbox}
\def\code#1{\texttt{#1}}


% insertar codigo en documento latex
%\usepackage{minted}
%\definecolor{LightGray}{gray}{0.92}
%\newminted{python}{ linenos,breaklines,mathescape,texcomments,xleftmargin=\parindent, numbersep=8pt,  bgcolor=LightGray
%}
%\renewcommand\listingscaption{C\'odigo}


% bibliografia: descomente estas dos lineas para usar estilo numerico, e.g. [1].
\usepackage[square,numbers,sort&compress]{natbib}
\bibliographystyle{apastyle}

\pagestyle{fancy}
\renewcommand{\footrulewidth}{0.4pt}
\renewcommand{\headrulewidth}{0.4pt}
\fancyfoot{}

%header y footer
\fancyfoot[RE,RO]{\thepage}
\fancyfoot[LE,LO]{Benjamin Ayala Baeza}


\definecolor{tcc}{RGB}{217,217,217} % Table cell color

\renewcommand\tabularxcolumn[1]{m{#1}}
\setlength{\arrayrulewidth}{0.5pt}
\renewcommand{\arraystretch}{2}

\tcbset{
    portada/.style={
        enhanced,
        colback=white,
        colframe=black,
        sharp corners,
        boxrule=0.4pt,
        drop shadow, % sombra elegante
        width=\textwidth,
        center
    }
}

%esto hace que desaperezca el header en la primera plana
\fancypagestyle{plainfirst}{
    \fancyhf{}
    \fancyfoot[RE,RO]{\thepage} % mantiene número de página
}

\begin{document}

\thispagestyle{plainfirst}

%\vspace*{\fill}
\begin{tcolorbox}[portada]
    \centering
    {\Huge \bfseries Métodos Númericos para la Física}\\[0.5cm]
    {\LARGE Derivación Númerica}\\[0.5cm]
    {\large Benjamín Ayala Baeza}
\end{tcolorbox}
%\vspace*{\fill}



\section{Porqué una de Derivación Númerica}

En la física uno de los recursos matematicos más utlizados es la derivada. Matematicamente la definimos como la taza de cambio de una función
y le tenemos una definición clásica, ánalitica y familiar. 

  \begin{align}
    \frac{d f(x)}{dx} = \lim_{h \rightarrow 0} \frac{f(x+h)- f(x)}{h} \label{eq:derivada-analitica}
  \end{align}

El problema ocurre cuando tenemos una situación en la que necesitamos generar una derivada donde tenemos un set de $n$ elementos 
discretos y si quisieramos estudiar el comportamiento de estos datos, podria ser util conocer su derivada, para esto aparece el termino
de derivación númerica la cual, en terminos mas simples, se basa en utilizar la expansión en serie de Taylor en un punto en
especifico para poder evaluar la derivada. 

\section{Tipos de Derivaciones Númericas}
Veremos que al utilizar la ecuación\eqref{eq:derivada-analitica} podemos manipularla para buscar la derivada en un punto especifico:

  \begin{align*}
    \left. \frac{d f(x)}{dx} \right|_{\tilde{x}} = \lim_{h \rightarrow 0} \frac{f(\tilde{x}+h)- f(\tilde{x})}{h} 
  \end{align*}

Intuitivamente tomariamos el valor de $h$ más pequeño, esto nos trae diversas problematicas, 
¿hay algún grado de error asociado?, ¿qué tan pequeño es suficiente?. Estas preguntas nos hacen darnos cuenta
que el approach anterior de simplemente evaluar en el punto con un valor muy pequeño no es lo más simple. \\
Por ello tendremos otros acercamientos a lo planteado incialmente, las series de Taylor.

  \subsection{Derivada Adelantada}
Haciendo no distinción entre $x$ y $\tilde{x}$, tomamos la expansión en serie de Taylor tal que:

\begin{align}
  f(x+h) = f(x) + h f'(x) + \frac{h^2}{2}f''(x) + \frac{h^3}{6} f'''(x) + \frac{h^4}{24}f^4(x) + \dots 
\end{align}

Que puede reescribirse como:

  \begin{align}
    f'(x) = \frac{f(x+h)-f(x)}{h} - \frac{h}{2}f''(x) + \dots
  \end{align}

Entonces, a partir de esto, podemos escribir la definición de la primera derivada adelantada:

\begin{tcolorbox}[colback=gray!20,colframe=black!40,title=Primera Derivada Adelantada]
  \begin{align}
    f' (x) = \frac{f (x+h)-f (x)}{h} - \mathcal{O} (h)
  \end{align}
\end{tcolorbox}

Este segundo termino del lado derecho, $\mathcal{O}(h)$ corresponde al error, el cual no deberia ser
muy grande o malo para un $h$ suficentemente pequeño. \\ Se llama derivada adelantada porque 
comienza en el valor $x$ y se mueve hacia adelante en la posicion positiva $x+h$.

\subsection{Derivada Retrasada}
Tal como en la sección anterior pero de forma inversa podemos definir lo que es la derivada retrasada,
utilizando la expansión en serie de Taylor para $f(x-h)$. La cual se puede definir tal como en la ecuación
de su contraparte como:

  \begin{align}
    f(x-h) = f(x) - h f'(x) + \frac{h^2}{2}f''(x) - \frac{h^3}{6} f'''(x) + \frac{h^4}{24}f^4(x) + \dots 
  \end{align}

Luego, obtendriamos:

  \begin{align}
    f'(x) = \frac{f(x)-f(x-h)}{h} + \frac{h}{2}f''(x) - \dots
  \end{align}

Luego, la aproximación de la derivada de $f(x)$ evaluada en $x$ sería:

\begin{tcolorbox}[colback=gray!20,colframe=black!40,title=Primera Derivada Retrasada]
  \begin{align}
    f' (x) = \frac{f (x)-f (x-h)}{h} + \mathcal{O} (h)
  \end{align}
\end{tcolorbox}

\subsection{Derivada Centrada}

Por último, buscamos un método que nos entregue una derivada más precisa que las anteriores, por ende 
al igual que en los casos anteriores buscamos la expansión en serie de Taylor de la funcion alrededor de $x$
pero esta vez buscamos un salto de valor $h/2$ en vez de solo $h$, tal que:

\begin{align}
  f \left( x + \frac{h}{2} \right) = f (x) + \frac{h}{2} f'(x) +\frac{h^2}{8}f''(x) + \frac{h^3}{48}f'''(x) + \dots \label{eq:derivada-centrada-positiva}
\end{align}

Podemos escribir la misma expansión pero moviendonos en la dirección negativa, tal que:

\begin{align}
  f \left( x - \frac{h}{2} \right) = f (x) - \frac{h}{2} f'(x) +\frac{h^2}{8}f''(x) - \frac{h^3}{48}f'''(x) + \dots \label{eq:derivada-centrada-negativa}
\end{align}

Si restamos la ecuación \eqref{eq:derivada-centrada-negativa} a la ecuación \eqref{eq:derivada-centrada-positiva} obtenemos:

\begin{tcolorbox}[colback=gray!20,colframe=black!40,title=Primera Derivada Centrada]
  \begin{align}
    f' (x) = \frac{f (x + \frac{h}{2})-f (x- \frac{h}{2})}{h} + \mathcal{O} (h^2)
  \end{align}
\end{tcolorbox}

Es útil notar que el uso de la derivada centrada corresponde a utilizar los puntos medios a los lados del valor donde estamos evaluando
y por ende si necesitamos evaluar la derivada o calcular en set finito, en algún endpoint no será posible porque es necesario conocer el siguiente valor a evaluar $n+1$.



\end{document}